% docs/manual/user/chapters/10-troubleshooting.tex
% 第 10 章：常见问题与排错

\chapter{常见问题与排错}

本章按"症状 → 诊断 → 修复"格式整理用户最常遇到的问题。如果你的具体错误消息没在这里出现，附录 E 有完整的错误消息索引。

\section{安装类}

\subsection{\texttt{cartopy} 编译失败}

\textbf{症状}：\cli{pip install openbench} 报 \texttt{Failed building wheel for cartopy}。

\textbf{诊断}：\modname{cartopy} 需要系统级 GEOS / PROJ 库。

\textbf{修复}：

\begin{minted}{bash}
# macOS
brew install geos proj
pip install openbench

# Ubuntu/Debian
sudo apt install libgeos-dev libproj-dev proj-data proj-bin
pip install openbench

# Conda 路径（绕过编译）
conda install -c conda-forge cartopy
pip install openbench
\end{minted}

\subsection{\texttt{PySide6} 装不上（HPC）}

\textbf{症状}：在 HPC 节点 \cli{pip install "openbench[gui]"} 报缺少图形库。

\textbf{修复}：HPC 上不要装 \cli{[gui]}，GUI 是本地客户端：

\begin{minted}{bash}
pip install openbench           # core only
pip install "openbench[remote]" # 如需远程 SSH
\end{minted}

\subsection{\texttt{ImportError: openbench[gui] not installed}}

\textbf{症状}：执行 \cli{openbench gui} 报 ImportError。

\textbf{修复}：装上 GUI extra：

\begin{minted}{bash}
pip install "openbench[gui]"
\end{minted}

\section{配置类}

\subsection{\texttt{Missing required section}}

\textbf{症状}：

\begin{minted}{text}
ConfigError: Missing required section: 'project'
\end{minted}

\textbf{修复}：四个必填顶层块（\yamlkey{project} / \yamlkey{evaluation} / \yamlkey{reference} / \yamlkey{simulation}）至少要有。详见第 3 章。

\subsection{\texttt{years must be a list of [start\_year, end\_year]}}

\textbf{症状}：YAML 写成 \yamlkey{years: 2010} 或 \yamlkey{years: [2010, 2011, 2012, 2013, 2014]}。

\textbf{修复}：必须恰好两个整数：

\begin{minted}{yaml}
project:
  years: [2010, 2014]   # ✓
\end{minted}

\subsection{\texttt{Variable 'X' has no entry in 'reference'}}

\textbf{症状}：\yamlkey{evaluation.variables} 列了变量，\yamlkey{reference} 没映射。

\textbf{修复}：

\begin{minted}{yaml}
evaluation:
  variables:
    - Latent_Heat        # 这个

reference:
  Latent_Heat: ERA5LAND  # 必须有这一行
\end{minted}

\subsection{\texttt{reference.sources must be a string}}

\textbf{症状}：写成嵌套形式 \yamlkey{reference.sources.<var>: ...}。

\textbf{修复}：第 3 章详解；reference 是 flat var→source 映射：

\begin{minted}{yaml}
# ✗
reference:
  sources:
    Evapotranspiration: GLEAM_v4.2a

# ✓
reference:
  Evapotranspiration: GLEAM_v4.2a
\end{minted}

\section{数据类}

\subsection{\texttt{Reference dataset not found}}

\textbf{症状}：\cli{check} 报 \texttt{not in registry}。

\textbf{诊断}：

\begin{minted}{bash}
# 看 catalog 里实际有哪些
openbench data list | grep -i gleam

# 检查 _LowRes / _MidRes 后缀解析
openbench data list --variable Evapotranspiration
\end{minted}

\textbf{修复}：用完整名（带 \texttt{\_LowRes} / \texttt{\_MidRes} 后缀）或调整 \yamlkey{project.grid_res} 让自动解析命中。

\subsection{Sim NC 文件找不到}

\textbf{症状}：run 时报 \texttt{No NC files matching pattern in /data/...}。

\textbf{诊断}：

\begin{minted}{bash}
# 实际的 NC 文件名是什么
ls /data/colm2024_run/ | head

# profile 期望的文件名模式
openbench model show CoLM2024
\end{minted}

\textbf{修复}：调整 \yamlkey{simulation.<label>.prefix} / \yamlkey{suffix}，或在 \yamlkey{variables} 里覆盖文件名模式。

\subsection{时间维度异常}

\textbf{症状}：评估突然 NaN 或维度错配。

\textbf{诊断}：\cli{ncdump -h <file.nc>} 看时间坐标的属性（\texttt{units}、\texttt{calendar}），与 catalog 中声明对照。

\textbf{常见根因}：

\begin{itemize}
  \item \texttt{calendar=noleap} vs \texttt{standard} 的混用
  \item \texttt{time} 单位是 \texttt{days since 1850-01-01} vs \texttt{seconds since 1970-01-01}
  \item \openbench{} 用 \modname{xarray} + \modname{cftime} 处理大多数情况；极端情况下需在 simulation entry 里手动指定时间偏移
\end{itemize}

\section{运行类}

\subsection{Out of memory}

\textbf{症状}：进程被 OOM killer 杀掉，\cli{run} 中途退出无清晰错误。

\textbf{诊断}：

\begin{minted}{bash}
# Linux
dmesg | grep -i "out of memory" | tail -5

# 看进程峰值
/usr/bin/time -v openbench run openbench.yaml
\end{minted}

\textbf{修复}：

\begin{itemize}
  \item 减小 \yamlkey{project.num_cores}（每 worker 占内存独立）
  \item 缩小 \yamlkey{lat_range} / \yamlkey{lon_range}（区域评估）
  \item 缩小 \yamlkey{years} 时段
  \item 用 \texttt{tim\_res: Month} 而非 \texttt{Day} / \texttt{Hour}（数据量降两个数量级）
\end{itemize}

\subsection{HDF5 lock error}

\textbf{症状}：

\begin{minted}{text}
OSError: Unable to open file (unable to lock file, errno = 11,
error message = 'Resource temporarily unavailable')
\end{minted}

\textbf{诊断}：HDF5 ≥ 1.14 默认对只读文件加锁；多 worker 同时读同一 reference NC 触发。

\textbf{修复}：\openbench{} 在 \cli{run} 启动时已设 \texttt{HDF5\_USE\_FILE\_LOCKING=FALSE}。如果你在脚本/Notebook 里手动调用 \modname{run\_evaluation}，自己设这个环境变量：

\begin{minted}{python}
import os
os.environ["HDF5_USE_FILE_LOCKING"] = "FALSE"
from openbench.runner.local import run_evaluation
\end{minted}

\subsection{运行卡住不动}

\textbf{诊断步骤}：

\begin{enumerate}
  \item 看 \file{output/<run>/run.log} 最后 30 行，确认是 \texttt{[INFO]} 卡住还是 \texttt{[DEBUG]} 还在动
  \item 用 \texttt{py-spy dump --pid \$(pgrep -f 'openbench run')} 看 Python 调用栈
  \item 看 CPU 与 IO：\texttt{top} / \texttt{iostat}。CPU = 0 + IO 不动 → 可能 NFS 挂掉；CPU 100\% + 一直跑 → 计算密集，等待
\end{enumerate}

\textbf{真正死锁较少见}，常见是大型 task 的真实计算时间被低估。

\subsection{时间对齐失败}

\textbf{症状}：

\begin{minted}{text}
TimeAlignmentError: No common time window between sim and ref
\end{minted}

\textbf{诊断}：

\begin{minted}{bash}
# Sim 的时段
ncdump -h /data/sim/file.nc | grep time

# Ref 的时段
ncdump -h /data/ref/file.nc | grep time
\end{minted}

\textbf{修复}：

\begin{itemize}
  \item 改 \yamlkey{project.years} 让窗口落在两者共有时段
  \item 切换 \yamlkey{time_alignment: per_pair} 让每对各算自己重叠
  \item 检查 \yamlkey{min_year_threshold}，如果重叠不足该值会被跳过
\end{itemize}

\section{站点类}

\subsection{站点匹配为空}

\textbf{症状}：站点评估报 \texttt{No matching stations}。

\textbf{诊断}：

\begin{minted}{bash}
# 查看站点列表
head /data/PLUMBER2/list/PLUMBER2.csv

# 站点经纬度是否落在 sim 的 lat/lon range
\end{minted}

\textbf{修复}：

\begin{itemize}
  \item 扩大 \yamlkey{lat_range} / \yamlkey{lon_range}
  \item 检查 fulllist CSV 的列名（必须有 \texttt{site, lat, lon}）
  \item 检查站点 NC 文件是否完整（缺站点会被跳过）
\end{itemize}

\subsection{HDF5 lock 在站点评估更明显}

站点评估 IO 模式是"打开同一文件 N 次拿不同站点切片"，比 grid 评估更容易触发 HDF5 lock。如果只在站点评估时报 lock 错误，几乎肯定是 \texttt{HDF5\_USE\_FILE\_LOCKING} 没设上。

\section{远程类}

\subsection{\cli{--remote} 立刻退出}

\textbf{症状}：

\begin{minted}{text}
Remote execution not yet implemented.
Install openbench[remote] and use openbench gui for remote execution.
\end{minted}

\textbf{当前状态}：v3.0a1 的 \cli{openbench run --remote} CLI 入口尚未实现。远程执行通过 GUI 提交。详见运维卷第 3 章。

\subsection{SSH 连接失败}

\textbf{症状}：GUI 中"Test connection"红色失败。

\textbf{诊断}：

\begin{minted}{bash}
# 在终端用同样的密钥手动 SSH
ssh -i ~/.ssh/my_key user@hpc01.example.edu
\end{minted}

\textbf{常见根因}：

\begin{itemize}
  \item 密钥权限不对（\texttt{chmod 600 \textasciitilde/.ssh/my\_key}）
  \item 主机 fingerprint 没接受过（先手动 SSH 一次）
  \item 防火墙 / VPN 隔离
\end{itemize}

\section{GUI 类}

\subsection{GUI 不启动}

第 8 章已详解。要点：

\begin{itemize}
  \item 装了 \cli{[gui]} extra
  \item 不在远程 SSH 无 X11 环境
  \item Linux 上设 \texttt{QT\_QPA\_PLATFORM=xcb}
\end{itemize}

\subsection{GUI 卡顿 / 无响应}

\textbf{诊断}：

\begin{itemize}
  \item Run Monitor 卡住通常是后端计算的事，不是 GUI 卡；等几分钟看进度条
  \item 配置页面卡顿可能是 registry 加载慢，看终端日志
\end{itemize}

\section{一般诊断流程}

遇到没在以上列出的错误，按以下顺序：

\subsection{Step 1：开 debug\_mode}

\begin{minted}{yaml}
project:
  debug_mode: true
\end{minted}

或临时：

\begin{minted}{bash}
PYTHONUNBUFFERED=1 openbench run openbench.yaml 2>&1 | tee debug.log
\end{minted}

\subsection{Step 2：看完整 log}

\begin{minted}{bash}
less +G output/<run>/run.log     # 跳到末尾
grep -i "error\|warn" debug.log
\end{minted}

\subsection{Step 3：缩小复现范围}

把 \yamlkey{evaluation.variables} 砍到 1 个、\yamlkey{years} 缩到 1 年、\yamlkey{simulation} 砍到 1 个。

\subsection{Step 4：dry-run + dump-config}

\begin{minted}{bash}
openbench run openbench.yaml --dry-run --dump-config
ls output/<run>/debug/
\end{minted}

\file{output/<run>/debug/} 下的中间 namelist 会暴露 v3.0 schema → legacy 翻译过程中的细节。

\subsection{Step 5：报告 Issue}

如果以上都没解决：

\begin{enumerate}
  \item 收集：\file{openbench.yaml}、\file{run.log}（关键段）、\cli{openbench version}、\cli{python --version}、OS
  \item 在 \href{https://github.com/CoLM-SYSU/OpenBench/issues}{GitHub Issues} 提交，附上以上信息
  \item 标注是 v3.0 用户、是否从 v2.0 迁移
\end{enumerate}

\section{完结}

到这里你已经完整走完用户卷的教程部分。后面的 5 个附录是参考性内容（配置项、CLI、数据集、模型、错误消息），按需查阅即可。

如果要更深入理解 \openbench{} 内部、添加新功能、或运维 HPC 部署，请参看：

\begin{itemize}
  \item \textbf{开发卷}（Developer Guide）：架构、扩展点、贡献流程
  \item \textbf{运维卷}（Operations Guide）：HPC 部署、远程 SSH、性能调优
\end{itemize}
